\chapter{Fazit und Ausblick}
\label{chap:fazit-ausblick}

\section{Zusammenfassendes Fazit}
\label{sec:fazit}
Die vorliegende Arbeit hat demonstriert, dass moderne cloud-native Betriebskonzepte auf Basis von Infrastructure as Code (HashiCorp Terraform) und GitOps (Argo~CD) auch unter strengen Budgetrestriktionen praxisnah und wissenschaftlich fundiert realisiert werden können. Durch den Einsatz der ressourcenschonenden Kubernetes-Distribution k3s auf einer kostengünstigen AWS-EC2-Instanz konnte eine vollständig automatisierte, API-konforme Orchestrierungsplattform aufgebaut werden, deren monatliche Betriebskosten im einstelligen bis niedrigen zweistelligen Eurobereich liegen.

Die Durchführung dreier kontrollierter Störungsexperimente lieferte eindeutige empirische Belege für die Resilienzeigenschaften des Gesamtsystems:
\begin{itemize}[noitemsep,topsep=2pt]
    \item Bei abrupten Prozessabstürzen sichert Kubernetes den raschen Wiederanlauf bei einer Verfügbarkeit von 97,9\,\% bis 98,8\,\% und einer effektiven Ausfallzeit von lediglich 0,5 Sekunden für die Basiskonfiguration.
    \item Defensiv konfigurierte Readiness-Probes und Rolling-Update-Strategien schützen den Cluster vor Totalausfällen bei fehlerhaften Releases, während ein Git-Revert den Ausgangszustand vollautomatisch wiederherstellt.
    \item Konfigurationsabweichungen durch unautorisierte Eingriffe im Cluster werden von Argo~CD in 3{,}85 Sekunden detektiert und autonom binnen 13{,}32 Sekunden korrigiert — weit unterhalb der Anforderungsgrenze FA-3 ($<$\,60\,s).
\end{itemize}

Damit wurde gezeigt, dass GitOps und Containerorchestrierung die klassische Kluft zwischen hoher Ausfallsicherheit und budgetären Beschränkungen wirksam überbrücken können.

\section{Ausblick und künftige Forschungsfragen}
\label{sec:ausblick}
Aufbauend auf den Ergebnissen dieser Untersuchung eröffnen sich vielversprechende Perspektiven für weiterführende Arbeiten:
\begin{itemize}[noitemsep,topsep=2pt]
    \item \textbf{Stochastisches Chaos Engineering:} Erweiterung der Testsuite um spezialisierte Frameworks wie Chaos Mesh zur Simulation zufälliger Netzwerklatenzen und CPU-Drosselungen im laufenden Betrieb.
    \item \textbf{Zustandsbehaftete Workloads:} Untersuchung von Failover-Mechanismen bei verteilten Datenbanken und Storage-CSI-Treibern (z.\,B. Longhorn oder Rook-Ceph) unter k3s.
    \item \textbf{Multi-Region- und Hybrid-Cloud-Szenarien:} Erprobung von GitOps-Workflows zur synchronen Verwaltung heterogener Cluster über mehrere Cloud-Provider oder On-Premises-Edge-Knoten hinweg.
\end{itemize}
